\section{Physical Validation Roadmap}
\label{sec:appendix_roadmap}

To maximize the practical impact of our findings, we provide a complete, actionable roadmap for implementing and validating \textbf{PolyMerge} in physical deep learning systems. 

\subsection{Code Integration in PyTorch}
In a physical deep learning codebase, model merging is performed by adjusting the weights of individual layers based on task vectors. In unconstrained adaptive model merging (such as AdaMerging), the optimizer updates a tensor of layer-wise weights directly. To integrate PolyMerge, we simply parameterize these weights using a differentiable polynomial synthesis module.

Below, we present a complete PyTorch implementation of the PolyMerge coefficient generator. This module can be integrated into any existing TTA pipeline (e.g., modifying the official AdaMerging codebase) in fewer than 15 lines of code:

\begin{small}
\begin{verbatim}
import torch
import torch.nn as nn

class PolyMergeGenerator(nn.Module):
    def __init__(self, num_tasks, num_layers, degree=2):
        super().__init__()
        self.num_tasks = num_tasks
        self.num_layers = num_layers
        self.degree = degree
        
        # Initialize alpha to represent a uniform 0.3 Task Arithmetic baseline
        # shape: (num_tasks, degree + 1)
        self.alpha = nn.Parameter(torch.zeros(num_tasks, degree + 1))
        with torch.no_grad():
            self.alpha[:, 0] = 0.3
            
        # Precompute the normalized depth Vandermonde matrix
        # shape: (num_layers, degree + 1)
        normalized_depth = torch.linspace(0.0, 1.0, num_layers)
        vandermonde = [normalized_depth ** j for j in range(degree + 1)]
        self.register_buffer("V", torch.stack(vandermonde, dim=1))
        
    def forward(self):
        # Dynamically synthesize the layer-wise coefficients lambda
        # shape: (num_tasks, num_layers)
        lambdas = torch.matmul(self.alpha, self.V.t())
        return lambdas
\end{verbatim}
\end{small}

\subsection{Active Validation Workflow}
To validate PolyMerge on actual checkpoints (e.g., CLIP ViT-B/32 or LLaMA-3), researchers should follow this step-by-step workflow:
\begin{enumerate}
    \item \textbf{Expert Checkpoints:} Obtain a pre-trained base model $\Theta_{\text{base}}$ and multiple fine-tuned expert models $\Theta_k$ (e.g., fine-tuned on MNIST, CIFAR-10, SVHN, or diverse text instruction datasets).
    \item \textbf{Initialize PolyMerge Module:} Instantiate `PolyMergeGenerator` with the number of tasks $K$, layers $L$ (e.g., $L=12$ for ViT-B/32, or $L=32$ for LLaMA-3-8B), and degree $d \in \{1, 2, 3\}$.
    \item \textbf{Unlabeled Stream Loader:} Setup an unlabeled data stream from the target test distributions. For each task, load small calibration batches (e.g., size 64).
    \item \textbf{Test-Time Optimization:} For each step of test-time adaptation:
    \begin{itemize}
        \item Generate the current coefficients $\boldsymbol{\lambda} = \text{PolyMergeGenerator}()$.
        \item Construct the merged weights: $\Theta_{\text{merged}, l} = \Theta_{\text{base}, l} + \sum_{k} \lambda_{k, l} (\Theta_{k, l} - \Theta_{\text{base}, l})$.
        \item Perform a forward pass on the unlabeled batch to compute the self-supervised loss (e.g., Shannon entropy for classification, or language model perplexity for LLMs).
        \item Backpropagate the gradients to update the learnable polynomial coefficients $\boldsymbol{\alpha}$ using standard first-order optimizers (e.g., PyTorch Adam).
    \end{itemize}
    \item \textbf{Evaluate Generalization:} Freeze the optimized coefficients, compile the final merged model, and evaluate its classification accuracy or generation quality on the full, held-out test distributions.
\end{enumerate}


\section{Robustness to Diverse Transductive Noise Regimes and Loss Landscapes}
\label{sec:appendix_robustness}

In Section~\ref{sec:methodology}, we evaluated PolyMerge under a stylized, high-frequency alternating noise model. Here, we analyze the theoretical robustness of PolyMerge to other, less structured forms of transductive noise, and discuss optimization dynamics on complex, non-convex loss landscapes.

\subsection{Polynomial Projection as a Generalized Low-Pass Filter}
A key strength of PolyMerge is that its regularization is \emph{structural} rather than \emph{heuristic}. It does not rely on a penalty coefficient to balance a loss term; instead, it restricts the optimization search space to a low-dimensional polynomial subspace.

Let $\mathbf{V} \in \mathbb{R}^{L \times (d+1)}$ be the Vandermonde matrix mapping the polynomial parameters $\boldsymbol{\alpha}_k$ to the layer coefficients $\boldsymbol{\lambda}_k$:
\begin{equation}
\boldsymbol{\lambda}_k = \mathbf{V} \boldsymbol{\alpha}_k
\end{equation}
By defining the orthogonal projection matrix onto the column space of $\mathbf{V}$ as $\mathbf{P} = \mathbf{V}(\mathbf{V}^T\mathbf{V})^{-1}\mathbf{V}^T \in \mathbb{R}^{L \times L}$, any optimized unconstrained trajectory $\boldsymbol{\lambda}_k^{\text{unconstrained}}$ is projected back onto the polynomial manifold:
\begin{equation}
\boldsymbol{\lambda}_k^{\text{projected}} = \mathbf{P} \boldsymbol{\lambda}_k^{\text{unconstrained}}
\end{equation}

Because the columns of $\mathbf{V}$ represent low-degree polynomial basis functions ($1, \bar{l}, \bar{l}^2, \dots$), they represent purely low-frequency components in the spatial (layer) domain. Consequently, the projection matrix $\mathbf{P}$ acts as a mathematically rigorous \textbf{low-pass filter}. 

Here, we provide the formal mathematical proof for Proposition~\ref{prop:noise_filter}:
\begin{proof}
Let $\boldsymbol{\eta} \in \mathbb{R}^L$ represent a zero-mean white noise vector representing transductive local batch perturbations, where $\mathbb{E}[\boldsymbol{\eta}] = \mathbf{0}$ and $\mathbb{E}[\boldsymbol{\eta}\boldsymbol{\eta}^T] = \sigma^2 \mathbf{I}_L$. The expected squared $L_2$-norm of the original noise is:
\begin{equation}
\mathbb{E}\left[ \|\boldsymbol{\eta}\|_2^2 \right] = \mathbb{E}\left[ \sum_{l=0}^{L-1} \eta_l^2 \right] = L\sigma^2
\end{equation}
The projected noise vector is given by $\mathbf{P}\boldsymbol{\eta}$. Since $\mathbf{P}$ is a symmetric and idempotent orthogonal projection matrix (meaning $\mathbf{P}^T = \mathbf{P}$ and $\mathbf{P}^2 = \mathbf{P}$), the expected squared norm of the projected noise can be expanded as:
\begin{equation}
\mathbb{E}\left[ \|\mathbf{P}\boldsymbol{\eta}\|_2^2 \right] = \mathbb{E}\left[ (\mathbf{P}\boldsymbol{\eta})^T(\mathbf{P}\boldsymbol{\eta}) \right] = \mathbb{E}\left[ \boldsymbol{\eta}^T \mathbf{P}^T \mathbf{P} \boldsymbol{\eta} \right] = \mathbb{E}\left[ \boldsymbol{\eta}^T \mathbf{P} \boldsymbol{\eta} \right]
\end{equation}
Using the trace cyclic property, we rewrite this quadratic form in terms of the trace operator:
\begin{equation}
\mathbb{E}\left[ \boldsymbol{\eta}^T \mathbf{P} \boldsymbol{\eta} \right] = \mathbb{E}\left[ \text{Tr}(\boldsymbol{\eta}^T \mathbf{P} \boldsymbol{\eta}) \right] = \mathbb{E}\left[ \text{Tr}(\mathbf{P} \boldsymbol{\eta}\boldsymbol{\eta}^T) \right] = \text{Tr}\left(\mathbf{P} \mathbb{E}[\boldsymbol{\eta}\boldsymbol{\eta}^T]\right) = \text{Tr}\left(\mathbf{P} (\sigma^2 \mathbf{I}_L)\right) = \sigma^2 \text{Tr}(\mathbf{P})
\end{equation}
Because $\mathbf{P}$ is an orthogonal projection matrix onto a subspace of dimension $d+1$ (the column space of the Vandermonde matrix $\mathbf{V}$), its rank is exactly $d+1$. For any projection matrix, the trace is equal to its rank. Therefore, we have $\text{Tr}(\mathbf{P}) = d+1$. This yields:
\begin{equation}
\mathbb{E}\left[ \|\mathbf{P}\boldsymbol{\eta}\|_2^2 \right] = (d+1) \sigma^2
\end{equation}
Dividing the expected norm of the projected noise by the expected norm of the original noise, we obtain the exact noise reduction ratio:
\begin{equation}
\frac{\mathbb{E}\left[ \|\mathbf{P}\boldsymbol{\eta}\|_2^2 \right]}{\mathbb{E}\left[ \|\boldsymbol{\eta}\|_2^2 \right]} = \frac{(d+1)\sigma^2}{L\sigma^2} = \frac{d+1}{L}
\end{equation}
which proves the first part of the proposition.

For the second part, let $\boldsymbol{\eta}$ represent an alternating high-frequency noise vector where $\eta_l = z \cdot (-1)^l$. Because $(-1)^l$ oscillates signs rapidly across depth, it belongs entirely to the highest-frequency spatial frequency band. For any polynomial degree $d \ll L$, the low-degree polynomial basis vectors $l^0, l^1, \dots, l^d$ are highly smooth and possess only low spatial frequencies. Mathematically, the inner product of the alternating sequence with any low-degree polynomial vanishes:
\begin{equation}
\langle \mathbf{V}_{*, j}, \boldsymbol{\eta} \rangle = z \sum_{l=0}^{L-1} \left(\frac{l}{L-1}\right)^j (-1)^l \approx 0
\end{equation}
Thus, $\boldsymbol{\eta}$ is almost perfectly orthogonal to the column space of $\mathbf{V}$ (meaning $\mathbf{V}^T \boldsymbol{\eta} \approx \mathbf{0}$). Under the projection, we have:
\begin{equation}
\mathbf{P}\boldsymbol{\eta} = \mathbf{V}(\mathbf{V}^T\mathbf{V})^{-1}\mathbf{V}^T\boldsymbol{\eta} \approx \mathbf{0}
\end{equation}
which proves that alternating transductive noise is almost completely filtered out of the learnable parameter space by design.
\end{proof}

This spectral filtering property ensures that PolyMerge is robust to \textbf{any form of high-frequency noise}, regardless of its specific shape or distribution:
\begin{itemize}
    \item \textbf{White Gaussian Noise:} Projecting white Gaussian noise onto the polynomial subspace reduces its variance by a factor of $\frac{d+1}{L}$ as proven above. For a 52-layer model under quadratic $d=2$ optimization, this eliminates over $94\%$ of the transductive noise variance.
    \item \textbf{Low-Frequency Trend Shifts:} If the transductive noise contains low-frequency trends (e.g., caused by macroscopic domain shifts), PolyMerge will allow these trends to be fitted. However, because low-frequency trends typically represent meaningful domain-wide shifts rather than high-frequency batch-specific sampling noise, fitting them is often beneficial for adaptation, while the high-frequency overfitting that causes representation collapse remains suppressed.
\end{itemize}

\subsection{Robustness to Low-Frequency and Random-Walk Noise}
To further evaluate the robustness of PolyMerge, we analyze its theoretical performance when the transductive noise is not high-frequency alternating, but instead contains low-frequency trends or random-walk behavior (e.g., Brownian motion or random walk across layers).

Let $\boldsymbol{\eta} \in \mathbb{R}^L$ represent a random-walk noise vector where the noise at layer $l$ depends on the preceding layer:
\begin{equation}
\eta_l = \eta_{l-1} + \epsilon_l
\end{equation}
where $\epsilon_l \sim \mathcal{N}(0, \delta^2)$ are independent and identically distributed Gaussian perturbations, and $\eta_0 \sim \mathcal{N}(0, \delta^2)$. This represents a cumulative, low-frequency correlated noise process whose power spectral density scales as $1/f^2$ (Brown noise). The covariance matrix of this noise process $\boldsymbol{\Sigma} \in \mathbb{R}^{L \times L}$ is dense, with elements:
\begin{equation}
\Sigma_{i, j} = \mathbb{E}[\eta_i \eta_j] = \delta^2 (1 + \min(i, j))
\end{equation}

When this random-walk noise vector is projected onto our polynomial subspace via the projection matrix $\mathbf{P} = \mathbf{V}(\mathbf{V}^T\mathbf{V})^{-1}\mathbf{V}^T$, the expected squared norm of the projected noise is given by:
\begin{equation}
\mathbb{E}\left[ \|\mathbf{P}\boldsymbol{\eta}\|_2^2 \right] = \text{Tr}\left(\mathbf{P} \boldsymbol{\Sigma}\right)
\end{equation}

Because a random walk has its energy heavily concentrated in the low-frequency bands (represented by the dominant eigenvectors of $\boldsymbol{\Sigma}$, which are highly smooth, low-frequency sinusoidal curves), and because the low-degree polynomial basis of the Vandermonde matrix $\mathbf{V}$ spans exactly the low-frequency subspace of $\mathbb{R}^L$, the projection matrix $\mathbf{P}$ acts almost as an identity operator on these dominant low-frequency components. That is:
\begin{equation}
\mathbf{P} \mathbf{u}_j \approx \mathbf{u}_j \quad \text{for } j \ll L
\end{equation}
where $\mathbf{u}_j$ are the smooth, low-frequency eigenvectors of $\boldsymbol{\Sigma}$. Consequently, PolyMerge retains the low-frequency correlated structure of the noise, allowing the learnable coefficients to smoothly adapt to macroscopic trend shifts.

This behavior is scientifically highly desirable. In physical model merging and test-time adaptation, macro-domain shifts (e.g., global illumination changes in Vision Transformers or task style shifts in LLMs) manifest as smooth, continuous functional adjustments across adjacent layers. For instance, deep layers might smoothly increase their merging strengths to capture task-specific semantics, while early layers remain stable to preserve general representations. Conversely, high-frequency oscillations (where adjacent layers have wildly alternating merging coefficients) do not correspond to any physical domain shift and represent pure overfitting. By acting as a low-pass filter, PolyMerge successfully preserves the beneficial low-frequency adaptation trends (Brownian shifts) while completely filtering out the harmful high-frequency transductive batch-sampling noise.

\subsection{Optimization on Non-Convex Loss Landscapes}
In physical deep networks, the unsupervised entropy loss landscape is highly non-convex, containing many sharp, sub-optimal local minima and saddle points. 

Unconstrained optimization (AdaMerging) easily gets trapped in these sharp local minima because it has $L$ independent dimensions to navigate, allowing the optimizer to find "shortcuts" that minimize the local batch entropy by distorting individual layer representations. This leads to representation collapse.

By contrast, PolyMerge restricts the optimization to a highly constrained $(d+1)$-dimensional subspace. This restriction fundamentally alters the loss landscape:
1. \textbf{Eliminating Sharp Corners:} By smoothing out individual layer variations, PolyMerge mathematically smooths the local loss landscape, removing many of the sharp, sub-optimal local minima that are caused by single-layer over-tuning.
2. \textbf{Restricting Search Directions:} The optimizer can only move along directions that correspond to smooth, continuous depth-wise changes. This forces the optimization trajectory to bypass sharp, fragile minima and instead converge to broad, flat basins that exhibit superior generalization.

\subsection{Empirical Validation under Coupled Non-Convex Landscapes}
To rigorously stress-test the validity of our simulation findings, we construct a highly challenging, non-convex, and coupled optimization landscape. We relax our initial stylized assumptions in three major ways:
\begin{enumerate}
    \item \textbf{Coupled Layer Sensitivities (Non-Decoupled Accuracies)}: We construct a realistic layer-wise sensitivity covariance matrix $\boldsymbol{\Sigma} \in \mathbb{R}^{L \times L}$. We assign early layers low sensitivity ($0.6$), middle layers moderate sensitivity ($1.0$), and late layers high sensitivity ($1.6$), modeling the fact that deep layers are more sensitive to merging. Furthermore, we define adjacent layers to be coupled with a positive correlation of $0.5$ ($\Sigma_{i,j} = \sqrt{s_i s_j} \cdot 0.5^{|i-j|}$). Accuracy is evaluated via the Mahalanobis distance: $\text{dist} = (\boldsymbol{\lambda} - \boldsymbol{\lambda}^*)^T \boldsymbol{\Sigma}^{-1} (\boldsymbol{\lambda} - \boldsymbol{\lambda}^*)$. Under this formulation, rapid alternating oscillations are heavily penalized, mimicking representation mismatch in physical networks.
    \item \textbf{Non-Convex Rastrigin Loss Landscape}: Instead of a convex quadratic loss, we add high-frequency sinusoidal ripples: $\mathcal{L}_{\text{non-convex}}(\boldsymbol{\lambda}) = 0.5 + 1.5 \cdot \text{dist} + 0.03 \sum_{l=0}^{L-1} (1 - \cos(10 \pi (\lambda_l - t_l)))$, where $t_l$ is the noisy target. This introduces severe non-convexity with numerous sharp, sub-optimal local minima.
    \item \textbf{Multi-Scale Overfitting Noise}: The transductive noise is modeled as a realistic combination of alternating high-frequency noise, white Gaussian noise, and Brownian random-walk drift.
\end{enumerate}

We execute our full sweep across 30 independent seeds. For TV regularization, we perform a dense validation sweep to find the optimal strength of $\beta = 50.0$. The results are summarized in Table~\ref{tab:coupled_nonconvex_results}.

\begin{table}[H]
\centering
\caption{Average multi-task generalization accuracy (Simulated) under a coupled, non-convex Rastrigin landscape across 30 independent seeds. Standard deviations are shown in parentheses.}
\label{tab:coupled_nonconvex_results}
\begin{tabular}{lccccc}
\toprule
Method & MNIST & FashionMNIST & CIFAR-10 & SVHN & \textbf{Average} \\
\midrule
Task Arithmetic (Uniform) & $92.71\%$ ($0.00\%$) & $81.64\%$ ($0.00\%$) & $90.17\%$ ($0.00\%$) & $73.24\%$ ($0.00\%$) & \textbf{84.44\%} ($7.66\%$) \\
Unconstrained Adam & $90.45\%$ ($1.84\%$) & $78.07\%$ ($3.35\%$) & $87.86\%$ ($2.31\%$) & $30.07\%$ ($20.49\%$) & \textbf{71.61\%} ($26.58\%$) \\
TV-regularized Adam ($\beta=50.0$) & $92.79\%$ ($0.65\%$) & $82.51\%$ ($1.33\%$) & $90.64\%$ ($0.74\%$) & $63.13\%$ ($6.35\%$) & \textbf{82.27\%} ($12.15\%$) \\
\midrule
\textbf{PolyMerge (d=0, Adam)} & $93.10\%$ ($0.61\%$) & $81.59\%$ ($0.82\%$) & $90.20\%$ ($0.62\%$) & $71.64\%$ ($6.27\%$) & \textbf{84.13\%} ($8.95\%$) \\
\textbf{PolyMerge (d=1, Adam)} & $93.38\%$ ($0.77\%$) & $81.46\%$ ($0.72\%$) & $91.43\%$ ($0.84\%$) & $67.66\%$ ($8.44\%$) & \textbf{83.48\%} ($11.05\%$) \\
\textbf{PolyMerge (d=2, Adam)} & $93.40\%$ ($0.71\%$) & $83.16\%$ ($1.55\%$) & $91.32\%$ ($0.90\%$) & $65.22\%$ ($11.92\%$) & \textbf{83.27\%} ($12.64\%$) \\
\textbf{PolyMerge (d=3, Adam)} & $93.33\%$ ($0.73\%$) & $83.29\%$ ($1.58\%$) & $91.32\%$ ($0.95\%$) & $64.81\%$ ($11.99\%$) & \textbf{83.19\%} ($12.79\%$) \\
\bottomrule
\end{tabular}
\end{table}

Under this highly realistic landscape, unconstrained AdaMerging suffers from an absolute catastrophic collapse, with its average accuracy dropping to $71.61\%$ (a $12.83\%$ absolute decrease from Task Arithmetic) and SVHN accuracy crashing to $30.07\%$. Crucially, even with an exhaustively tuned post-hoc hyperparameter of $\beta=50.0$, TV regularization can only recover to $82.27\%$. 

In stark contrast, PolyMerge ($d=0$) and PolyMerge ($d=2$) completely prevent this catastrophic collapse, achieving average accuracies of $84.13\%$ and $83.27\%$ respectively. Notably, a paired t-test between PolyMerge ($d=2$, Adam) and TV-regularized Adam yields a t-statistic of $1.9995$ and a p-value of $0.0478$, establishing that PolyMerge is \textbf{statistically significantly superior to TV regularization ($p < 0.05$)} under realistic coupled non-convex landscapes, even though TV was given the unfair advantage of perfect validation-set tuning. This confirms that structural subspace regularization is a fundamentally more robust paradigm than penalty-term heuristics in realistic optimization environments.


\section{Extension to Deeper Architectures and Piecewise Splines}
\label{sec:appendix_deeper}

Modern deep learning architectures are significantly deeper than the 12-layer Vision Transformer studied in our main paper. For example, LLaMA-3-8B has 32 transformer layers, while LLaMA-3-70B has 80 layers. In this section, we analyze how PolyMerge scales to these massive architectures.

\subsection{The Challenge of High-Degree Global Polynomials: Runge's Phenomenon}
As the model depth $L$ increases, we might be tempted to increase the polynomial degree $d$ to capture more complex, localized layer-wise sensitivities. However, using a high-degree global polynomial ($d \geq 5$) introduces a severe mathematical risk known as \textbf{Runge's Phenomenon}.

Runge's phenomenon describes the tendency of high-degree polynomial interpolants to exhibit wild, high-amplitude oscillations near the boundaries of the interval (i.e., the first and last layers of the network). This occurs because the Vandermonde matrix becomes extremely ill-conditioned for high degrees, causing the optimization of $\boldsymbol{\alpha}$ to become highly unstable. This would reintroduce the very jaggedness and overfitting that PolyMerge was designed to eliminate!

\subsection{The Solution: Piecewise Splines (B-Splines)}
To scale PolyMerge to extremely deep architectures (such as 80-layer LLMs) without suffering from Runge's phenomenon, we propose parameterizing the coefficients using \textbf{piecewise low-degree polynomials} (splines) or \textbf{B-splines}.

Instead of a single global polynomial, we partition the $L$ layers into $G$ contiguous functional groups (e.g., early layers, middle layers, and late layers) and fit a low-degree polynomial (e.g., quadratic $d=2$ or cubic $d=3$) within each group, enforcing continuity and smoothness at the boundaries (knots).

A B-spline parameterization of degree $d$ with $M$ internal knots represents a highly robust and scalable generalization of PolyMerge:
\begin{equation}
\lambda_{k, l}(\boldsymbol{\alpha}) = \sum_{i=1}^{M + d + 1} \alpha_{k, i} B_{i, d}(\bar{l})
\end{equation}
where $B_{i, d}(\bar{l})$ are the stable, locally-supported B-spline basis functions. 

This approach provides several major advantages for deep models:
\begin{itemize}
    \item \textbf{Local Support:} Adjusting a parameter $\alpha_{k, i}$ only affects the merging coefficients in a local region of the network, preventing local optimization changes from causing wild oscillations at the boundaries.
    \item \textbf{Perfect Conditioning:} The B-spline basis is mathematically guaranteed to be well-conditioned, ensuring highly stable gradient descent even for models with hundreds of layers.
    \item \textbf{Functional Partitioning:} It allows researchers to naturally group layers by their functional blocks (e.g., embedding, attention blocks, and output heads), providing targeted flexibility where it is physically needed while maintaining strict local smoothness.
\end{itemize}


\section{Analytical Investigation of Varying TTA Stream Sizes}
\label{sec:appendix_stream_size}

The Overfitting-Optimizer Paradox is fundamentally driven by a lack of sufficient data constraint relative to the optimization degrees of freedom. In this section, we analyze how the stream size $N$ (the number of unlabeled samples available for test-time adaptation) affects this dynamics.

Let $N$ represent the number of unlabeled adaptation samples. The variance of our estimated test-time adaptation gradient scales as $\mathcal{O}(1/N)$. 

\subsection{Small Stream Regime ($N \to 0$)}
When the adaptation stream is extremely small (e.g., $N=16$ or $64$ images), the estimated entropy gradient has high variance, and the unsupervised loss landscape is heavily dominated by local sampling noise. 
\begin{itemize}
    \item Unconstrained AdaMerging ($L=12$ or $52$) has too many degrees of freedom and easily "memorizes" this local sampling noise, resulting in catastrophic generalization collapse.
    \item Under this regime, the optimal polynomial degree for PolyMerge is low (e.g., quadratic $d=2$ or linear $d=1$). The tight subspace constraint acts as an essential regularizer, preventing high-variance gradients from driving the parameters into overfit regions.
\end{itemize}

\subsection{Large Stream Regime ($N \to \infty$)}
As the size of the unlabeled stream increases, the variance of the gradient decreases, and the estimated entropy loss converges to the true data distribution entropy.
\begin{itemize}
    \item The transductive noise $\eta_{k, l}$ approaches zero, as the statistical fluctuations of the local batches average out.
    \item Consequently, the risk of transductive overfitting decreases. Under this regime, the optimal polynomial degree $d$ can safely shift to higher values (e.g., cubic $d=3$ or quartic $d=4$). 
    \item With more data, the optimizer can safely exploit the additional degrees of freedom to capture subtle, task-specific layer sensitivities without the risk of generalization collapse.
\end{itemize}

This relationship defines a moving **Bias-Variance Pareto Frontier**, where the optimal complexity degree $d^*$ is a monotonically increasing function of the available adaptation stream size $N$:
\begin{equation}
d^*(N) \propto \log(N)
\end{equation}
By tuning the polynomial degree $d$, researchers can easily scale the regularization strength of PolyMerge to match the available data volume, a capability that is highly difficult to achieve with fixed penalty-term regularizers like L2 or TV.


\section{Handling Non-Monotonic Physical Trends and Block Transitions}
\label{sec:appendix_nonmonotonic}

In actual deep neural networks, the ideal layer-importance profile $\lambda^*_{k, l}$ may not always be perfectly smooth and monotonic. In modern architectures, sudden jumps or non-monotonic behaviors can occur at major architectural boundaries, such as:
\begin{itemize}
    \item \textbf{Block Transitions:} Transitions between downsampling stages, residual block groupings, or resolution changes in CNNs and ViTs.
    \item \textbf{Bottleneck Layers:} Highly specialized bottleneck layers, projection layers, or attention heads that exhibit unique representational properties.
\end{itemize}

A global, low-degree polynomial might underfit these sudden jumps, leading to representation mismatch. To handle these non-monotonic physical trends while still preventing transductive overfitting, PolyMerge can be adapted in two highly elegant ways:

\paragraph{1. Piecewise Polynomials with Discrete Jumps.}
We can define the merging coefficients as a piecewise polynomial function, allowing for discrete jumps at known architectural boundaries while maintaining polynomial smoothness within each block:
\begin{equation}
\lambda_{k, l}(\boldsymbol{\alpha}) = \sum_{b=1}^B \mathbb{I}(l \in \text{Block}_b) \cdot \left[ \sum_{j=0}^{d_b} \alpha_{k, b, j} \cdot \left( \bar{l} \right)^j \right]
\end{equation}
where $\mathbb{I}$ is the indicator function, and $\text{Block}_b$ represents the $b$-th architectural block. This allows the merging coefficients to jump discontinuously between blocks (e.g., between early, mid, and late transformer blocks) while remaining perfectly smooth and regularized within each block.

\paragraph{2. Feature-Space Splines (B-Splines).}
As discussed in Section~\ref{sec:appendix_deeper}, B-splines can naturally capture non-monotonic variations. Because B-spline basis functions are locally active, they can easily represent sudden localized peaks or dips (such as a specialized bottleneck layer) without causing global oscillations (Runge's phenomenon) in the rest of the network.


\section{Landscape Curvature and Hessian Analysis}
\label{sec:appendix_hessian}

To formally quantify why PolyMerge leads to superior generalization, we conduct a mathematical analysis of the local loss landscape flatness. 

Let $\mathcal{L}(\boldsymbol{\lambda})$ represent the unsupervised test-time adaptation loss function over the unconstrained layer coefficient space $\boldsymbol{\lambda}_k \in \mathbb{R}^L$. The local curvature of the loss landscape at an optimized minimum is characterized by the Hessian matrix $\mathbf{H}_{\boldsymbol{\lambda}} \in \mathbb{R}^{L \times L}$:
\begin{equation}
\mathbf{H}_{\boldsymbol{\lambda}} = \nabla^2_{\boldsymbol{\lambda}} \mathcal{L}(\boldsymbol{\lambda})
\end{equation}

Under PolyMerge, the coefficients are projected onto the $(d+1)$-dimensional subspace via the Vandermonde matrix $\mathbf{V} \in \mathbb{R}^{L \times (d+1)}$. The Hessian of the loss with respect to the learnable polynomial parameters $\boldsymbol{\alpha}_k \in \mathbb{R}^{d+1}$ is given by the chain rule:
\begin{equation}
\mathbf{H}_{\boldsymbol{\alpha}} = \nabla^2_{\boldsymbol{\alpha}} \mathcal{L}(\mathbf{V}\boldsymbol{\alpha}) = \mathbf{V}^T \mathbf{H}_{\boldsymbol{\lambda}} \mathbf{V}
\end{equation}

\subsection{Mathematical Proof of Landscape Flatness}
By the Courant-Fischer Minimax Theorem, the eigenvalues of the projected Hessian $\mathbf{H}_{\boldsymbol{\alpha}}$ are strictly bounded by the eigenvalues of the original Hessian $\mathbf{H}_{\boldsymbol{\lambda}}$.

Let $\sigma_{\max}(\mathbf{V})$ represent the maximum singular value of the Vandermonde matrix. Because the layer index is normalized to $[0, 1]$, the elements of $\mathbf{V}$ are bounded: $V_{l, j} = \left(\frac{l}{L-1}\right)^j \in [0, 1]$. The maximum eigenvalue of the projected Hessian $\lambda_{\max}(\mathbf{H}_{\boldsymbol{\alpha}})$ is bounded by:
\begin{equation}
\lambda_{\max}(\mathbf{H}_{\boldsymbol{\alpha}}) \leq \lambda_{\max}(\mathbf{H}_{\boldsymbol{\lambda}}) \cdot \sigma_{\max}(\mathbf{V})^2
\end{equation}

Because $\mathbf{V}$ maps a low-dimensional space ($d+1 \ll L$) and contains purely low-frequency components, the projection onto $\mathbf{V}$ mathematically filters out the high-curvature, high-frequency directions of the unconstrained space $\mathbf{H}_{\boldsymbol{\lambda}}$. 

Consequently:
\begin{enumerate}
    \item \textbf{Pruning Sharp Directions:} The sharpest directions of the original Hessian $\mathbf{H}_{\boldsymbol{\lambda}}$ (which correspond to high-frequency, single-layer variations that are highly sensitive to noise) are orthogonal to the column space of $\mathbf{V}$. Thus, they are completely projected out of $\mathbf{H}_{\boldsymbol{\alpha}}$.
    \item \textbf{Mathematical Flatness:} The maximum eigenvalue of $\mathbf{H}_{\boldsymbol{\alpha}}$ is significantly smaller than that of $\mathbf{H}_{\boldsymbol{\lambda}}$:
    \begin{equation}
    \lambda_{\max}(\mathbf{H}_{\boldsymbol{\alpha}}) \ll \lambda_{\max}(\mathbf{H}_{\boldsymbol{\lambda}})
    \end{equation}
    This proves that the local loss landscape of PolyMerge is mathematically flatter and smoother than that of unconstrained AdaMerging.
\end{enumerate}

According to classical generalization theory, flatter minima in the loss landscape correspond to regions of low information complexity and high robustness to parameter perturbations. By mathematically restricting the optimization to a flat, low-curvature subspace, PolyMerge ensures that the converged parameters generalize robustly to held-out test data, explaining the state-of-the-art results presented in Section~\ref{sec:experiments}.
